\documentclass[a4paper,12pt]{article}
\newcounter{CJKLanguage}
\setcounter{CJKLanguage}{3}
\input{LaTeX-ToolKit/template.tex}
\begin{document}
\title{口語文法}
\author{沈威宇}
\date{\temtoday}
\titletocdoc
\sct{口語文法}
\ssc{Linguistics}
\sssc{Pronunciation}
\begin{itemize}
\item \ruby{訓読}{くんよ}み is the way of reading \ruby{漢字}{かんじ} characters using the native Japanese word that matches the meaning of the Chinese character when it was introduced.
\item \ruby{音読}{おんよ}み is the way of reading \ruby{漢字}{かんじ} characters based on the original Chinese pronunciation of the character.
\begin{itemize}
\item \ruby{呉音}{ごおん}: 音読み that are transmitted to Japanese from Chinese before 漢音.
\item \ruby{漢音}{かんおん}: 音読み that are transmitted to Japanese from Chinese during the 隋 and 唐 dynasties by 遣隋使・遣唐使.
\item \ruby{唐音}{とうおん}/\ruby{宋音}{そうおん}: 音読み that are transmitted to Japanese from Chinese since the 宋 dynasty.
\end{itemize}
\item \ruby{熟字訓}{じゅくじくん}/\ruby{義訓}{ぎくん}: Readings of kanji combinations that have no direct correspondence to the characters' individual 音読み and 訓読み, e.g., 昨日（きのう）、今日（きょう）、明日（あした）.
\end{itemize}
\sssc{語種}
日本語の\ruby{語種}{ごしゅ} includes:
\begin{itemize}
\item \ruby{和語}{わご}: Words origin in the Japanese language.
\item \ruby{漢語}{かんご}: Words borrowed from Chinese during earlier historical periods.
\item \ruby{外来語}{がいらいご}: Words borrowed from other languages excluding Chinese, frequently written using 片仮名.
\end{itemize}
\ssc{用言の活用}
\sssc{体言と用言}
\begin{itemize}
\item 体言は自立語で、活用がなく、主語となる語つまり名詞を主に指す。
\item 用言は自立語で、活用があり、述語や連体修飾語となる語つまり動詞・形容詞・形容動詞を主に指す。
\end{itemize}
\sssc{活用}
\begin{longtable}[c]{|C{0.09\textwidth}|C{0.09\textwidth}|C{0.09\textwidth}|C{0.09\textwidth}|C{0.09\textwidth}|C{0.09\textwidth}|C{0.09\textwidth}|C{0.08\textwidth}|C{0.09\textwidth}|}
\hline
 & 語幹 & 基本形/辞書形 & 未然形 & 連用形 & 終止形 & 連体形 & 仮定形 & 命令形 \\\hline\endhead
後続する主な言葉 & - & - & ない・ん・れる・られる・せる・させる・う・よう & ます・に・た・だ・て・で & 「。」・が・から・けれど・けれども・と・し & の・名詞 & ば &  \\\hline
だ & - & だ & だろ（+う） & だっ（+た）/で（+（係助詞）ある/as 中止形） & だ & な & なら & - \\\hline
です & - & です & でしょ（+う） & でし（+た） & です & です & なし & なし \\\hline
れる & - & れる & れ & れ & れる & れる & れれ & れろ・れよ \\\hline
られる & - & られる & られ & られ & られる & られる & られれ & られろ・られよ \\\hline
せる & - & せる & せ & せ & せる & せる & せれ & せろ・せよ \\\hline
させる & - & させる & させ & させ & させる & させる & させれ & させろ・させよ \\\hline
う & - & う & - & - & う & う & - & - \\\hline
よう & - & よう & - & - & よう & よう & - & - \\\hline
ます & - & ます & ませ（+ん）/ましょ（+う） & まし（+た・て） & ます & ます & ますれ & ませ/まし \\\hline
ない & - & ない & なかろ（+う） & なかっ（+た）・なく（+て） & ない & ない & なけれ & - \\\hline
め（ん） & - & め（ん） & - & ず & め（ん） & め（ん） & ね & - \\\hline
た（だ） & - & た（だ） & たろ（だろ） & - & た（だ） & た（だ） & たら（だら） & - \\\hline
形容詞/イ形容詞 & 母音で終わる & 語幹い & 語幹かろ（+う） & 語幹く & 語幹い & 語幹い & 語幹けれ & - \\\hline
\ruby{良}{い}い & \ruby{良}{い}い & \ruby{良}{い}い & \ruby{良}{よ}かろ & \ruby{良}{よ}く & \ruby{良}{い}い & \ruby{良}{い}い & \ruby{良}{よ}けれ & - \\\hline
形容動詞/ナ形容詞 & 母音で終わる & 語幹 & 語幹だろ（+う） & 語幹だっ（+た）/語幹で（+（係助詞）ある/as 中止形） & 語幹だ & 語幹な & 語幹なら & - \\\hline
五段活用動詞/グループ1動詞 except ワア行 & 子音で終わる & 語幹u & 語幹a（+ない・れる・せる）/語幹o（+う） & 語幹i（+ます・に, and for サ行五段, た・だ・て・で）/「ござる・なさる・くださる・いらっしゃる・おっしゃる」のます形は語幹-r+い（+ます）です/語幹-子音+い（カ、ガ行）・っ（\ruby{行}{い}く、タ、ラ行）・ん（ナ、バ、マ行）（+た・だ・て・で） & 語幹u & 語幹u & 語幹e & 語幹e \\\hline
ワア行五段活用動詞/グループ1動詞 & 母音で終わる & 語幹う & 語幹わ（+ない・れる・せる）/語幹お（+う） & 語幹い（+ます・に）/語幹っ（+た・だ・て・で） & 語幹う & 語幹う & 語幹え & 語幹え \\\hline
上、下一段活用動詞/グループ2動詞 & 母音（上一段活用動詞i, 下一段活用動詞e）で終わる & 語幹る & 語幹 & 語幹 & 語幹る & 語幹れ & 語幹れ & 語幹よ/ろ \\\hline
カ行変格活用動詞/グループ3動詞 & \ruby{来}{く}r & \ruby{来}{く}る & \ruby{来}{こ}（+ない・られる・させる・よう） & \ruby{来}{き}（+ます・に・た・だ・て・で） & \ruby{来}{く}る & \ruby{来}{く}る & \ruby{来}{く}れ & \ruby{来}{こ}い \\\hline
サ行変格活用動詞/グループ3動詞 & すr & する & し（+ない・られる・させる・よう） & し（+ます・に・た・だ・て・で） & する & する & すれ & しろ/せよ \\\hline
\end{longtable}
中止形: 連用形 or て形 used as a conjunction.
\sssc{助動詞}
\begin{itemize}
\item だ：断定する意を表す。
\item です：「だ」の丁寧語。
\item 「じゃ」は「では」（「だ」の連用形「で」+係助詞「は」）の音便化。「ではない」「じゃない」「ではありません」「じゃありません」は否定の意味を示す表現。「ではない」「じゃない」は普通形で、「で」+「は」+「ない」（動詞「ある」の否定形）から成り。「ではありません」「じゃありません」は丁寧形で、「で」+「は」+「ありません」（動詞「ある」の丁寧否定形）から成り。
\item れる・られる：受身・尊敬・自発・可能の意を表す。
\item せる・させる：使役の意を表す。
\item う・よう ：推量・意志・勧誘の意を表す。
\item ます：丁寧の意を表す。
\item ない：打ち消しの意を表す。
\item め（ん）：打ち消しの意を表す。
\item た・だ：過去・完了・存続・確認の意を表す。ガ・ナ・バ・マ行の五段活用動詞に続くときだけ「だ」になります。
\end{itemize}
\sssc{Last 仮名 of 辞書形}
\begin{itemize}
\item 辞書形 of 形容詞/イ形容詞 end in 「い」. However, there exist 形容動詞/ナ形容詞 whose 辞書形 end in 「い」, e.g., きれい、嫌い.
\item 辞書形 of 一段活用動詞 and 変格活用動詞 end in 「る」. However, there exist 五段活用動詞 whose 辞書形 end in 「る」, e.g., 作る、始まる、終わる、乗る、撮る、帰る、割る.
\end{itemize}
\ssc{助詞}
\sssc{格助詞}
名詞とそれ以外の単語との意味関係（格）を表す。
\begin{longtable}[c]{|C{0.15\textwidth}|C{0.35\textwidth}|C{0.35\textwidth}|}
\hline
助詞 & 働き & 例文 \\\hline\endhead
が & 主語 & 友達が来ました。 \\\hline
が & 自動詞の対象 & 寿司が食べます。 \\\hline
を & 他動詞の対象 & ジュースを買います。 \\\hline
と & 共同の相手 & 友人と食べる。 \\\hline
と & 引用 & 彼は「はい」と答えた。 \\\hline
に & 移動の着点 & 学校に行く。 \\\hline
に & 動作の相手 & 部下に任せる。 \\\hline
に & 授受の受け手（与え手が受け手に物を「あげる」系統/「くれる」系統） & 花子は太郎にプレゼントをあげた。 \\\hline
に & 授受の与え手（受け手が与え手に物を「もらう」系統） & 太郎は花子にプレゼントをもらう。 \\\hline
に & 時間（「〜ごろ」は「に」をつけてもつけなくてもいい） & 8時に起きる。 \\\hline
に & 到達点 & 学校に行く。 \\\hline
に & 目的 & 流れ星を見に行く。 \\\hline
へ （「え」と発音する） & 方向 & 東へ向かう。 \\\hline
へ （「え」と発音する） & 到達点。「に」と区別が曖昧だが、「に」よりも対象への強調があり、それが物であるときは使いにくい。 & 学校へ行く。 \\\hline
で & 場所 & 図書館で本を読みます。 \\\hline
で & 手段 & バスで行きます。 \\\hline
から & 空間的起点 & 台湾から東京まで飛行機で行きます。 \\\hline
から & 時間的起点 & 毎朝8時から学校へ行きます。 \\\hline
から & 経由する場所 & 東京を出て、名古屋から京都へと向かう。 \\\hline
から & 原因 & 操作ミスから事故が生じた。 \\\hline
から & 材料、構成要素 & 米から酒ができる。 \\\hline
\end{longtable}
\sssc{並列助詞}
2つ以上のものを並立させる。（格助詞に含める説もある）
\begin{longtable}[c]{|C{0.15\textwidth}|C{0.35\textwidth}|C{0.35\textwidth}|}
\hline
助詞 & 働き & 例文 \\\hline
と & 全て挙げる & コーヒーと紅茶と牛乳を飲む。 \\\hline
か & 選択（１つを選ぶ） & コーヒーか紅茶か牛乳を飲む。 \\\hline
や (also 、; など after last item is optional) & 全て挙げない & コーヒーや紅茶や牛乳（など）を飲む。 \\\hline
だの & 全て挙げない。二つ目は必要ではない。「…だの…だの」の形で用いられるが、「…だの…など」の形で用いられることもある。 & コーヒーだの紅茶だの牛乳だのを飲む。 \\\hline
\end{longtable}
\sssc{副助詞}
体言や副詞、格助詞の後などにつき全体として副詞的に働く。
\begin{longtable}[c]{|C{0.1\textwidth}|C{0.2\textwidth}|C{0.5\textwidth}|}
\hline
助詞 & 働き & 例文 \\\hline\endhead
まで（漢字表記も「迄」） & 空間的到達点 & 台湾から東京まで飛行機で行きます。 \\\hline
まで（漢字表記も「迄」） & 時間的な限度 & 15日までにご連絡ください。 \\\hline
など & 多くの中の一例を挙げる & コーヒーや紅茶や牛乳などがある。 \\\hline
くらい・ぐらい（漢字表記も「位」） & 大凡の分量や程度、基準、事態 & 学校まで30分くらいかかります。 \\\hline
\end{longtable}
\sssc{係助詞}
特定の文の要素を強調する。（副助詞に含める説もある）
\begin{longtable}[c]{|C{0.15\textwidth}|C{0.35\textwidth}|C{0.35\textwidth}|}
\hline
助詞 & 働き & 例文 \\\hline\endhead
は（「わ」と発音する） & 主題 & 私は学生です。 \\\hline
は（「わ」と発音する） & 対比 & コーヒーは飲みますが、紅茶\textbf{は}飲みません。 \\\hline
も & 同様のことを並べる。また、全てを列挙しない場合は暗に同様のことがあることを示す働きもする。 & 父も母も不在です。 \\\hline
も & 極端なものを提示し、それ以外のことについても当然成り立つことを暗示する。 & そんなこともできるのか。 \\\hline
も & 最も程度の極端な場合の想定を表す。 & 最低でも一次リーグは突破したい。 \\\hline
も & 肯定文中で数量の多さ・程度の高さを強調する。 & 鈴木さん家の家賃15万もするらしいよ。 \\\hline
も & 否定文中で数量が少ないこと・程度が低いことを強調する。 & １０メートルも泳げない。 \\\hline
も & 不定称の語について全面的な否定及び肯定などを示す。 & そのことは誰もが知っている。 \\\hline
\end{longtable}
\sssc{接続助詞}
文と文の意味関係を表して接続するもの。
\begin{longtable}[c]{|C{0.15\textwidth}|C{0.35\textwidth}|C{0.35\textwidth}|}
\hline
助詞 & 働き & 例文 \\\hline\endhead
ば & 仮定の順接 & 雨が降れば、行きません。 \\\hline
ば & 一般条件 & 春になれば、桜が咲きます。 \\\hline
ば & 並立の関係 & 雨が降れば、風が吹く。 \\\hline
と & 仮定の順接 & 雨が降ると、行きません。 \\\hline
と & 一般条件 & 春になると、桜が咲きます。 \\\hline
と & 仮定の逆接 & 雨が降ると、行きます。 \\\hline
て（ガ・ナ・バ・マ行の五段活用動詞に続くときだけ「で」になります。） & 確定の順接 & 雨が降って、行きません。 \\\hline
て（ガ・ナ・バ・マ行の五段活用動詞に続くときだけ「で」になります。） & 動作や作用の継続 & 学校に行って勉強する。 \\\hline
て（ガ・ナ・バ・マ行の五段活用動詞に続くときだけ「で」になります。） & 並立の関係 & 雨が降って、風が吹く。 \\\hline
て（ガ・ナ・バ・マ行の五段活用動詞に続くときだけ「で」になります。） & 補助の関係 & 嫌になってしまう。 \\\hline
ので & 確定の順接 & 雨が降っているので、行きません。 \\\hline
のに & 確定の逆接 & 雨が降っているのに、行きます。 \\\hline
から & 確定の順接 & 雨が降るから、行きません。 \\\hline
ても（ガ・ナ・バ・マ行の五段活用動詞に続くときだけ「でも」になります。） & 仮定の逆接 & 雨が降っても、行きます。 \\\hline
ても（ガ・ナ・バ・マ行の五段活用動詞に続くときだけ「でも」になります。） & 確定の逆接 & 雨が降っても、行きます。 \\\hline
が & 確定の逆接 & 雨が降るが、行きます。 \\\hline
が & 並立の関係 & 雨が降るが、風が吹く。 \\\hline
けれど・けれども & 確定の逆接 & 雨が降るけれど（も）、行きます。 \\\hline
けれど・けれども & 並立の関係 & 雨が降るけれど（も）、風が吹く。 \\\hline
し & 並立の関係 & 雨が降るし、風が吹く。 \\\hline
\end{longtable}
\sssc{終助詞}
文の最後にきて話し手の感情や態度を表す。
\begin{longtable}[c]{|C{0.15\textwidth}|C{0.35\textwidth}|C{0.35\textwidth}|}
\hline
助詞 & 働き & 例文 \\\hline\endhead
か（上昇調） & 疑問 & 今、何時ですか。 \\\hline
か（上昇調） & 反語 & なぜあなたを忘れようか。 \\\hline
か（下降調） & 相槌 & そうですか。 \\\hline
か/よ/ね & 感動 & 合格したか/よ/ね。おめでとう。 \\\hline
ね & 確認を求める & 今日は暑いですね。 \\\hline
ね & 答えとして確認を表す & 今日は暑いですね。そうだ\textbf{ね}。 \\\hline
よね & 確認を求める（確信がもてない） & この映画、おもしろかったよね。 \\\hline
よ/ね & 告知 & 10時からですよ/ね。 \\\hline
よ & 勧誘 & いっしょに行こうよ。 \\\hline
よ & 命令 & 危ないよ。 \\\hline
\end{longtable}
\ssc{数詞}
\sssc{漢語系}
0 to 10:
\begin{itemize}
\item 零（れい、ぜろ）、一（いち）、二（に）、三（さん）、四（し、よん）、五（ご）、六（ろく）、七（しち、なな）、八（はち）、九（く、きゅう）、十（じゅう）。
\item 零（ぜろ）は外来語, 四（よん）、七（なな）は和語, and the others are 漢語. However, this set of numbers is called 漢語系.
\item 四 is typically pronounced よん instead of し; 七 is typically pronounced なな instead of しち.
\end{itemize}
Attachments of 一：
\begin{itemize}
\item Positive integer powers of 万 begin with 一 if no digit would otherwise precede.
\item If 千 directly precedes the name of positive integer powers of 万, 一 is attached before 千, which yields 一千（いっせん）.
\item If 千 does not directly precede the name of positive integer powers of 万, attaching 一 is optional.
\item 一 is normally not attached before 百 or 十.
\end{itemize}
Phonetic modifications:
\begin{itemize}
\item 促音便: いち, はち, じゅう before the か, さ, た, or は series, and ろく before the か and は series, undergo 促音便 to いっ, はっ, じゅっ, and ろっ.
\item 半濁音化 and 濁音化 after っ: After 促音便, the は series is shifted to ぱ seires, and sometimes the か series is shifted to が series.
\item 半濁音化 and 濁音化 after ん: After さん, なん, せん, まん, (and, less commonly, よん), 清音 are shifted to 濁音, or, for は series, 半濁音.
\item よん is sometimes modified to よ.
\end{itemize}
Hundred to trillion:
\begin{itemize}
\item 百（ひゃく）: 促音便, 半濁音化 after っ, 濁音化 after ん.
\item 千（せん）: 促音便, 濁音化 after ん.
\item 万（まん）
\item 億（おく）
\item 兆（ちょう）: 促音便.
\end{itemize}
Strings of digits:
\begin{itemize}
\item Lengthen normally unimoraic digits such as に, ご, and less used, し and く to form groups of two digits of exactly four morae as long as possible, e.g., 五四二六二 is pronounced ごうよんにいろくに.
\item - between digits is pronounced の.
\item 零 is also read as まる and written as 〇.
\end{itemize}
\sssc{和語系}
一（ひと）、二（ふた）、三（み）、四（よ）、五（いつ）、六（む）、七（なな）、八（や）、九（ここの）、十（とお）。
\sssc{大字}
大字（だいじ）: Only 壱 (1), 弐 (2), 参 (3), and 拾 (10) are stipulated by Japanese law.
\sssc{助数詞}
助数詞 is used after 数詞 or 何（なん）(or 幾（いく）).

In principle,
\begin{itemize}
\item 数詞 before 漢語や外来語の助数詞 are in 漢語系.
\item 1 to 10 (sometimes 2, sometimes 4) before 和語の助数詞 are in 和語系, often with phonetic modifications.
\item Some 助数詞 has two forms that are used with 数詞 of 和語系 and 漢語系 respectively and for fixed ranges.
\end{itemize}
But this is with plenty exceptions.

A list of 助数詞: 漢語系 with よん and なな and 何 unless otherwise specified.
\begin{itemize}
\item つ: 一つ（ひとつ）、二つ（ふたつ）、三つ（みっつ）、四つ（よっつ）、五つ（いつつ）、六つ（むっつ）、七つ（ななつ）、八つ（やっつ）、九つ（ここのつ）、十（とお）. General-purpose 助数詞 for 1 to 10 for objects or things when a more specific counter is unnecessary or unknown. Cannot be used for numbers more than 10. いくつ.
\item 歳/才（さい）: 1 to 10 can also be called using つ; 何歳/才 also and more respectfully おいくつ; 20 years old is more commonly called 二十歳/二十（はたち）.
\item 人: 一人（ひとり）、二人（ふたり）for 1 and 2; 漢語 人（にん）for others.
\item 秒（びょう） and 秒間（びょうかん）
\item 分（ふん） and 分間（ふんかん）: 促音便; 半濁音化 after っ; 半濁音化 after ん including よん. 〜時半 means 〜時三十分.
\item 時（じ） and 時間（じかん）: れい、よ、しち (時間 also なな)、く. Both 12 hr and 24 hr systems are used. For the former, 午前 and 午後 are added before 〜時 for a.m. and p.m. respectively.
\item 日（にち） and 日間（にちかん）: 一日（ついたち）for the date and 一日（いちにち）for one day. 二日（ふつか）、三日（みっか）、〜四日（よっか）、五日（いつか）、六日（むいか）、七日（なのか）、八日（ようか）、九日（ここのか）、十日（とおか）、二十日（はつか）for 2 to 10, 20, and those end with 4; 漢語 日（にち）for others, with 〜十七日（じゅうしちにち）、〜十九日（じゅうくにち）; as day also 日間, with 間（かん） added after. 一昨々日（さきおととい）$\rightarrow$ 一昨日（おととい）$\rightarrow$ 昨日（きのう）$\rightarrow$ 今日（きょう）$\rightarrow$ 明日（あした）$\rightarrow$ 明後日（あさって）$\rightarrow$ 明々後日（しあさって）.
\item 月（がづ）(name of month): し、しち、く.
\item か月（かげつ）(also ヶ月（かげつ）・箇月（かげつ）・カ月（かげつ）, and less used, ヵ月（かげつ）・ケ月（かげつ）) (month): 促音便. 先々月（せんせんげつ）$\rightarrow$ 先月（せんげつ）$\rightarrow$ 今月（こんげつ）$\rightarrow$ 来月（らいげつ）$\rightarrow$ 再来月（さらいげつ）.
\item 年（ねん） and 年間（ねんかん）: よ、しち (年間 also なな). 西暦（せいれき）(Common Era), 紀元前（きげんぜん）(Before Common Era), and 元号（げんごう）such as 令和（れいわ）can be added before 〜年. The first year of a 元号 is called 元年（がんねん）. 〜年 without prefix means Common Era. 一昨々年（さきおととし）$\rightarrow$ 一昨年（おととし）$\rightarrow$ 去年（きょねん）$\rightarrow$ 今年（ことし）$\rightarrow$ 来年（らいねん）$\rightarrow$ 再来年（さらいねん）$\rightarrow$ 明明後年（みょうみょうごねん）.
\item 週間（しゅうかん）: 促音便. 先々週（せんせんしゅう）$\rightarrow$ 先週（せんしゅう）$\rightarrow$ 今週（こんしゅう）$\rightarrow$ 来週（らいしゅう）$\rightarrow$ 再来週（さらいしゅう）.
\item 階（かい）(floor (American)): 促音便; 濁音化 after ん. 地下（ちか） can be added before 〜階 for underground.
\end{itemize}
\ssc{コソアド}
\begin{longtable}[c]{|C{0.16\textwidth}|C{0.16\textwidth}|C{0.16\textwidth}|C{0.16\textwidth}|C{0.16\textwidth}|}
\hline
 & 話し手の近く & 聞き手の近く & それ以外 & 疑問語 \\\hline
連体詞名詞修飾 & この & その & あの & どの \\\hline
連体詞属性 & こんな & そんな & あんな & どんな \\\hline
丁寧形指示代名詞事物 & こちら & そちら & あちら & どちら（も3つ以上の選択） \\\hline
普通形指示代名詞事物 & これ & それ & あれ & どれ（も2つの選択） \\\hline
丁寧形指示代名詞人 & こちらの方/こちら/この方 & そちらの方/そちら/その方 & あちらの方/あちら/あの方 & どちらの方/どちら/どの方 \\\hline
普通形指示代名詞人 & この人/こいつ & その人/そいつ & あの人/あいつ & どの人/どいつ \\\hline
丁寧形指示代名詞場所 & こちら & そちら & あちら & どちら \\\hline
普通形指示代名詞場所 & ここ & そこ & あそこ & どこ \\\hline
丁寧形指示代名詞方向 & こちら & そちら & あちら & どちら \\\hline
普通形指示代名詞方向 & こっち & そっち & あっち & どっち \\\hline
副詞 & こう & そう & あう & どう \\\hline
\end{longtable}
\end{document}
